\documentclass[11pt,a4paper]{article}
\usepackage[utf8]{inputenc}
\usepackage[margin=1in]{geometry}
\usepackage{amsmath,amssymb}
\usepackage{graphicx}
\usepackage{booktabs}
\usepackage{listings}
\usepackage{xcolor}
\usepackage{float}
\usepackage{caption}
\usepackage{hyperref}
\usepackage{pbox}

% --- Custom Colors for Source Code Listing ---
\definecolor{mGreen}{rgb}{0,0.6,0}
\definecolor{mGray}{rgb}{0.5,0.5,0.5}
\definecolor{mPurple}{rgb}{0.58,0,0.82}
\definecolor{backgroundFix}{rgb}{0.96,0.96,0.96}

\lstdefinestyle{CStyle}{
    backgroundcolor=\color{backgroundFix},
    commentstyle=\color{mGreen},
    keywordstyle=\color{blue},
    numberstyle=\tiny\color{mGray},
    stringstyle=\color{mPurple},
    basicstyle=\ttfamily\footnotesize,
    breakatwhitespace=false,         
    breaklines=true,                 
    captionpos=b,                    
    keepspaces=true,                 
    numbers=left,                    
    numbersep=5pt,                  
    showspaces=false,                
    showstringspaces=false,
    showtabs=false,                  
    tabsize=2,
    language=C
}

\title{
    \textbf{Laporan Tugas Besar 1 IF4021 Pemodelan dan Simulasi} \\
    \vspace{0.2cm}
    \large Pemodelan Simulasi Sistem Car Rental Bus Menggunakan SIMLIB dan C
}

\author{
    \textbf{Brian A. Hadian} \\ 13523048 \\
    \textit{Kelompok 8}
    \makeatother
}
\date{\today}

\begin{document}

\maketitle

\hrule
\vspace{0.5cm}

% ==============================================================================
% DELIVERABLE 1: DESKRIPSI MASALAH (HALAMAN 1)
% ==============================================================================
\section*{1. Deskripsi dan Interpretasi Masalah}

Sistem pengangkut penumpang pada area penyewaan mobil (\textit{car-rental system}) menghubungkan tiga lokasi utama: Terminal 1 (Lokasi 1), Terminal 2 (Lokasi 2), dan Car Rental (Lokasi 3). Lintasan berbentuk sirkuit tertutup sejauh 10 mil merupakan rute untuk 1 unit bus berkapasitas maksimal 20 orang dengan kecepatan konstan 30 mil/jam (0.5 mil/menit). Bus bergerak searah jarum jam balik (\textit{counterclockwise}) melintasi rute $3 \to 1 \to 2 \to 3$.

\subsection*{Karakteristik Kedatangan dan Destinasi Penumpang}
\begin{itemize}
    \item Lokasi 1 (Terminal 1): Tingkat kedatangan penumpang berdistribusi eksponensial dengan laju $\lambda_1 = 14$ orang/jam. Seluruh penumpang bertujuan ke Car Rental (Lokasi 3).
    \item Lokasi 2 (Terminal 2): Tingkat kedatangan penumpang berdistribusi eksponensial dengan laju $\lambda_2 = 10$ orang/jam. Seluruh penumpang bertujuan ke Car Rental (Lokasi 3).
    \item Lokasi 3 (Car Rental): Tingkat kedatangan penumpang berdistribusi eksponensial dengan laju $\lambda_3 = 24$ orang/jam. Destinasi dipilih acak: Terminal 1 dengan peluang $0.583$ dan Terminal 2 dengan peluang $0.417$.
\end{itemize}

\subsection*{Aturan Operasional Bus dan Pelayanan}
\begin{enumerate}
    \item Penurunan (\textit{Unloading}): Penumpang diturunkan secara FIFO. Waktu penurunan per orang berdistribusi Uniform(16, 24) detik (Stream 4).
    \item Penaikan (\textit{Loading}): Penumpang dinaikkan secara FIFO hingga kapasitas bus terpenuhi (maksimum 20 orang). Waktu naik per orang terdestribusi Uniform(15, 25) detik (Stream 5).
    \item Aturan Waktu Dwell Minimum: Bus wajib berhenti minimal 5 menit (300 detik) di setiap lokasi. Jika dalam 5 menit tidak ada proses naik/turun yang berlangsung, bus akan berangkat.
    \item Waktu Tempuh Perjalanan:
    \begin{itemize}
        \item Lokasi 3 ke Lokasi 1: $4.5\text{ mil} / 0.5\text{ mil/min} = 9.0\text{ menit} = 540\text{ detik}$.
        \item Lokasi 1 ke Lokasi 2: $1.0\text{ mil} / 0.5\text{ mil/min} = 2.0\text{ menit} = 120\text{ detik}$.
        \item Lokasi 2 ke Lokasi 3: $4.5\text{ mil} / 0.5\text{ mil/min} = 9.0\text{ menit} = 540\text{ detik}$.
    \end{itemize}
\end{enumerate}

Simulasi dijalankan selama 80 jam ($288,000$ detik) untuk mengumpulkan statistik antrean, waktu tunda, okupansi bus, waktu henti bus, waktu \textit{loop}, dan total waktu penumpang dalam sistem.

\newpage

% ==============================================================================
% PLACEHOLDERS GAMBAR (FIGURE PLACEHOLDERS)
% ==============================================================================

\section*{2. Diagram Sistem dan Tangkapan Layar Program}

\begin{figure}[H]
    \centering
    % SPACE FOR IMAGE: Figure 2.72 System Diagram
    \vspace{0.1cm}
    \begin{center}
        \includegraphics[width=0.55\textwidth]{fig2_72.png}
    \end{center}
    \vspace{0.1cm}
    \caption{Diagram Arsitektur Sistem Car-Rental (Fig. 2.72)}
    \label{fig:car_rental_diagram}
\end{figure}

\vspace{0.2cm}

\begin{figure}[H]
    \centering
    % SPACE FOR IMAGE: Program Execution Screenshot
    \vspace{0.1cm}
    \begin{center}
        \includegraphics[width=0.65\textwidth]{execution_snapshot1.png}
    \end{center}
    \vspace{0.1cm}
    \caption{Tangkapan Layar Hasil Eksekusi Program Simulasi C (SIMLIB) 1}
    \label{fig:execution_screenshot}
\end{figure}

\vspace{0.2cm}

\begin{figure}[H]
    \centering
    % SPACE FOR IMAGE: Program Execution Screenshot
    \vspace{0.1cm}
    \begin{center}
        \includegraphics[width=0.65\textwidth]{execution_snapshot2.png}
    \end{center}
    \vspace{0.1cm}
    \caption{Tangkapan Layar Hasil Eksekusi Program Simulasi C (SIMLIB) 2}
    \label{fig:execution_screenshot}
\end{figure}

\newpage

% ==============================================================================
% DELIVERABLE 2: SOURCE CODE (SIMLIB C)
% ==============================================================================
\section*{3. Kode Sumber Program (\texttt{car\_rental.c})}

Berikut adalah kode sumber lengkap program simulasi sistem car-rental dalam bahasa C menggunakan pustaka SIMLIB.

\lstinputlisting[style=CStyle, caption={Kode Sumber Lengkap Simulation Engine (\texttt{car\_rental.c})}]{car_rental.c}

\newpage

% ==============================================================================
% DELIVERABLE 3: OUTPUT PROGRAM & HASIL STATISTIK
% ==============================================================================
\section*{4. Output Program dan Pembahasan Statistik}

Hasil eksekusi simulasi selama 80 jam ($288,000$ detik) disajikan pada tabel-tabel statistik berikut:

\subsection*{(a) Statistik Jumlah Penumpang dalam Antrean (\textit{Queue Length})}
\begin{table}[H]
\centering
\begin{tabular}{lccc}
\toprule
\textbf{Lokasi / Stasiun} & \textbf{Rata-rata (Orang)} & \textbf{Maksimum (Orang)} & \textbf{Minimum (Orang)} \\
\midrule
Lokasi 1 (Terminal 1) & 7.4457 & 27.0 & 0.0 \\
Lokasi 2 (Terminal 2) & 4.9424 & 17.0 & 0.0 \\
Lokasi 3 (Car Rental) & 9.4496 & 30.0 & 0.0 \\
\bottomrule
\end{tabular}
\caption{Statistik Jumlah Penumpang Antre di Setiap Lokasi}
\end{table}

\subsection*{(b) Statistik Waktu Tunda Antrean (\textit{Queue Delay})}
\begin{table}[H]
\centering
\begin{tabular}{lcccc}
\toprule
\textbf{Lokasi} & \textbf{Rata-rata (s)} & \textbf{Maksimum (s)} & \textbf{Minimum (s)} & \textbf{Jumlah Sampel} \\
\midrule
Lokasi 1 & 1904.89 (31m 45s) & 5534.83 (92m 15s) & 18.60 & 1120 \\
Lokasi 2 & 1809.34 (30m 09s) & 7876.40 (131m 16s) & 15.31 & 784 \\
Lokasi 3 & 1467.80 (24m 28s) & 4073.23 (67m 53s) & 22.71 & 1854 \\
\bottomrule
\end{tabular}
\caption{Statistik Waktu Tunda Penumpang Menunggu Bus}
\end{table}

\subsection*{(c) Statistik Jumlah Penumpang di Dalam Bus (\textit{Bus Occupancy})}
\begin{table}[H]
\centering
\begin{tabular}{lccc}
\toprule
\textbf{Parameter} & \textbf{Rata-rata (Orang)} & \textbf{Maksimum (Orang)} & \textbf{Kapasitas Maksimal} \\
\midrule
Okupansi Bus & 13.8530 & 20.0 & 20 \\
\bottomrule
\end{tabular}
\caption{Statistik Penggunaan Kapasitas Tempat Duduk Bus}
\end{table}

\subsection*{(d) Statistik Waktu Henti Bus di Setiap Lokasi (\textit{Bus Dwell Time})}
\begin{table}[H]
\centering
\begin{tabular}{lcccc}
\toprule
\textbf{Lokasi} & \textbf{Rata-rata (s)} & \textbf{Maksimum (s)} & \textbf{Minimum (s)} & \textbf{Jumlah Berhenti} \\
\midrule
Lokasi 1 & 401.89 (6m 42s) & 610.24 (10m 10s) & 300.00 (5m 00s) & 110 \\
Lokasi 2 & 328.39 (5m 28s) & 557.20 (9m 17s) & 300.00 (5m 00s) & 110 \\
Lokasi 3 & 684.11 (11m 24s) & 824.59 (13m 45s) & 334.31 (5m 34s) & 110 \\
\bottomrule
\end{tabular}
\caption{Statistik Durasi Bus Berhenti (\textit{Dwell Time}) Per Stasiun}
\end{table}

\subsection*{(e) Statistik Waktu Putaran Bus (\textit{Loop Time})}
\begin{table}[H]
\centering
\begin{tabular}{lcccc}
\toprule
\textbf{Metrik} & \textbf{Rata-rata (s)} & \textbf{Maksimum (s)} & \textbf{Minimum (s)} & \textbf{Total Putaran} \\
\midrule
Waktu Loop ($3 \to 3$) & 2614.40 (43m 34s) & 2884.15 (48m 04s) & 2142.88 (35m 43s) & 110 \\
\bottomrule
\end{tabular}
\caption{Statistik Waktu Tempuh Satu Putaran Lengkap Bus}
\end{table}

\subsection*{(f) Statistik Total Waktu Penumpang Dalam Sistem (\textit{Total Time in System})}
\begin{table}[H]
\centering
\begin{tabular}{lcccc}
\toprule
\textbf{Lokasi Asal} & \textbf{Rata-rata (s)} & \textbf{Maksimum (s)} & \textbf{Minimum (s)} & \textbf{Jumlah Penumpang} \\
\midrule
Lokasi 1 & 3115.64 (51m 56s) & 6737.62 (112m 18s) & 1064.31 (17m 44s) & 1120 \\
Lokasi 2 & 2732.66 (45m 33s) & 8820.25 (147m 00s) & 835.71 (13m 56s) & 784 \\
Lokasi 3 & 2503.43 (41m 43s) & 5297.35 (88m 17s) & 691.89 (11m 32s) & 1838 \\
\bottomrule
\end{tabular}
\caption{Statistik Total Waktu Penumpang Sejak Datang Hingga Tiba di Tujuan}
\end{table}

\end{document}